\documentclass[12pt,a4paper]{article}
\usepackage[margin=2.5cm]{geometry}
\usepackage{amsmath,amssymb,booktabs,longtable,hyperref,xcolor,enumitem,microtype}
\usepackage{graphicx,caption,subcaption}
\hypersetup{colorlinks=true,linkcolor=blue,citecolor=blue,urlcolor=blue}
\setlength{\parskip}{0.5em}
\setlength{\parindent}{0pt}

% ─────────────────────────────────────────────────────────────────────────────
% STATUS LABEL COMMANDS
% ─────────────────────────────────────────────────────────────────────────────
\newcommand{\blocked}{\textsc{blocked}}
\newcommand{\conditional}{\textsc{conditional}}
\newcommand{\derivedcheck}{\textsc{derived\_check}}
\newcommand{\unknown}{\textsc{unknown}}
\newcommand{\assumed}{\textbf{[ASSUMED]}}
\newcommand{\designspec}{\textbf{[DESIGN\_SPECIFIED]}}
\newcommand{\notinstalled}{\textbf{[NOT\_INSTALLED]}}
\newcommand{\mfrspec}{\textbf{[MANUFACTURER\_SPEC]}}
\newcommand{\litbound}{\textbf{[LITERATURE\_BOUND]}}
\newcommand{\unverified}{\textbf{[UNVERIFIED]}}
\newcommand{\superseded}{\textbf{[SUPERSEDED~v3.0]}}

\newcommand{\verdictbox}[1]{%
  \vspace{0.5em}\par
  \noindent\colorbox{yellow!30}{\parbox{\dimexpr\linewidth-2\fboxsep}{\bfseries #1}}\par
  \vspace{0.5em}
}

\title{\textbf{Quantum Thermal Architecture: A Mode-Separated Validation Framework\\
for Cryogenic LCVD and NV-Based Helium Isotope Sensing}}
\author{[Author list]}
\date{Manuscript version v4.0 --- \today}

\begin{document}
\maketitle

% ─────────────────────────────────────────────────────────────────────────────
\verdictbox{GLOBAL STATUS:
Current: BLOCKED (required hardware not installed; required measurements not performed).
Post-install forecast: CONDITIONAL ($\tau_c \geq 292\,\mu$s AND $C_\mathrm{contr}>0$ at 10\,mK required).
Validated system: NOT AVAILABLE.
Breakthrough claim: NOT MADE.
Full-cycle MC (current): 0.0\%.
PASS count: 0.}
% ─────────────────────────────────────────────────────────────────────────────

% =============================================================================
\begin{abstract}
% =============================================================================
This work proposes a same-chamber, mode-separated cryogenic validation framework
combining an LCVD-capable process mode (Mode~B) with NV-center-based helium
isotope sensing (Mode~D)~\cite{Pobell2007,DegenRMP2017}.
\textbf{The framework is not experimentally validated.}
The current state is \blocked{} because required hardware is not installed
and required measurements are missing.

LCVD growth and NV sensing are mutually exclusive, hardware-interlocked
operating modes (enforced by IL-01, IL-02, IL-14; see Section on interlocks).
They do not operate simultaneously.
The simulation is an assumption-tracked forecast model, not a first-principles proof.
No gate reaches \textsc{Pass} without verified in-system measurement.

Detection feasibility is jointly gated by two co-equal scientific unknowns:
\begin{enumerate}[noitemsep]
  \item $\tau_c$ of adsorbed $^3$He on the actual F-terminated diamond surface
        at 10\,mK \cite{Bernu2006,MorhardPRB1996} [\unknown{}; canonical threshold $\tau_c \geq 292\,\mu$s].
  \item $C_\mathrm{contr}$ at 10\,mK under actual NV readout conditions
        \cite{DegenRMP2017,Sangtawesin2019} [\unknown{}].
\end{enumerate}
\textbf{If $C_\mathrm{contr} = 0$ or too low at 10\,mK, detection fails
regardless of $\tau_c$.}

Unresolved engineering bottlenecks include: $G_\mathrm{eff}$ of the
Ag-sinter/diamond interface \assumed{}, residual gas suppression
(bakeout not executed) \notinstalled{}, cryotrap \designspec{},
NEG pump \designspec{}, optical absorption $\eta_\mathrm{abs}$ \assumed{},
microwave dissipation $P_\mathrm{mw}$ \assumed{},
vibration PSD at NV position $S_\mathrm{vib}$ \mfrspec{},
NV depth $d_\mathrm{NV}$ \assumed{}, $T_2^*$ \assumed{},
surface spin density $n_s$ \assumed{}, and mode interlock behavior \designspec{}.

Gate counts: 0~\textsc{Pass} $|$ 47~\conditional{} $|$ 0~\textsc{Fail}
$|$ 2~\unknown{} $|$ 23~\blocked{} $|$ 11~\derivedcheck{} (D9 gas kinetics; 10 non-lumped multiphysics numerical checks).
\medskip

\noindent\textbf{Canonical gate-table summary (single source of truth, generated by
\texttt{qta\_full\_sim.py}):} 83 total gates --- 0 PASS, 47 CONDITIONAL, 23 BLOCKED,
2 UNKNOWN, 11 DERIVED\_CHECK. Canonical $\tau_c$ threshold: 292\,$\mu$s (v3.3). The
27.728\,$\mu$s value (v3.0) is SUPERSEDED, NOT\_CANONICAL, NOT\_LIVE\_GATE\_LOGIC.

\end{abstract}

% =============================================================================
\section{Scope and Non-Claims}
\label{sec:nonclaims}
% =============================================================================

This manuscript does \textbf{not} claim:
\begin{itemize}[noitemsep]
  \item working millikelvin LCVD hardware
  \item demonstrated NV helium spin-noise detection
  \item mode-switched LCVD and NV sensing
  \item validated in-situ growth at 10\,mK
  \item verified $\tau_c$ for $^3$He on the actual diamond surface
  \item verified $C_\mathrm{contr}$ at 10\,mK
  \item verified Ag-sinter $G_\mathrm{eff}$
  \item verified residual gas suppression
  \item verified QCM coverage tracking
  \item verified optical absorption or transient recovery
  \item a breakthrough, Nobel-worthy result, or DARPA-ready system
  \item a validated engineering system
  \item a proof of feasibility
\end{itemize}

\textbf{Correct category:} blocked conditional validation framework.

The framework defines a mode-separated architecture and a program of
experiments required to test whether the sensing premise survives contact with
physical reality.

% =============================================================================
\section{Canonical Mode Map}
\label{sec:modes}
% =============================================================================

The architecture enforces strict temporal separation:

\begin{center}
\begin{tabular}{llp{8cm}}
\toprule
Mode & Label & Description \\
\midrule
A & BASELINE/STABILIZATION & Cryogenic baseline at the operating set-point; thermal/vacuum/NV-baseline verification; readiness for any exposure or processing. Sensing OFF. LCVD OFF. \\
B & MATERIAL\_PROCESSING/LCVD\_GROWTH & LCVD on; CH$_4$ or other precursor present; fs-laser processing; pulsed molecular-beam delivery; sensing off; heat switch open. Helium ABSENT. \\
C & ISOLATION\_PURGE\_THERMAL\_RECOVERY & Growth inputs off; gas lines isolated; cryopump/baffle residual species; sample thermally relaxes; vibration settles; RGA/QCM/witness-coupon checks. \\
D & SENSING/MEASUREMENT & LCVD off; precursor below threshold; He-3 dosed; NV ODMR / Ramsey at the millikelvin sensing condition. \\
\bottomrule
\end{tabular}
\end{center}

\textbf{LCVD growth and NV sensing are not simultaneous.
They occur as mutually exclusive, hardware-interlocked modes inside the same
cryogenic chamber.}
Mode~B and Mode~D cannot run together.
Mode~C exists to isolate contamination and allow thermal recovery before
Mode~D begins; Mode~A is the cryogenic baseline/stabilization state.

All mode transitions and interlocks are \designspec{} unless physically
implemented and tested.

% =============================================================================
\section{Evidence Status Vocabulary}
\label{sec:vocab}
% =============================================================================

\begin{description}
  \item[\blocked{}] Required hardware, measurement, or validation infrastructure missing; result cannot be interpreted.
  \item[\conditional{}] Model forecasts feasibility only if unverified assumptions later measure correctly.
  \item[\unknown{}] Parameter is too unknown to bound; no reliable forecast possible.
  \item[\derivedcheck{}] Internal mathematical or gas-kinetic check is self-consistent; not a physical system pass.
  \item[\textsc{Fail}] Modeled or measured value violates threshold.
  \item[\textsc{Pass}] Allowed \emph{only} after verified in-system measurement under relevant conditions.
\end{description}

\textbf{Absolute rule:} No gate reaches \textsc{Pass} from \assumed{},
\unknown{}, \litbound{}, \mfrspec{}, \designspec{}, or \notinstalled{} inputs.

% =============================================================================
\section{Committed Mechanical and Thermal Architecture}
\label{sec:support}
% =============================================================================

\subsection{Rejected Baseline: Direct Vespel SP-22 Supports}

Direct Vespel SP-22 polyimide supports are \textbf{rejected} and are not part
of the committed architecture.
They create an unacceptable direct conductive path from 4\,K to the 10\,mK
sample node.
All Vespel SP-22 calculations appear in Appendix~\ref{app:vespel_rejected}
for reference only.

\subsection{Committed Support Architecture}

\textbf{The committed sample-support architecture replaces direct Vespel rods
with a hybrid Kevlar~49 / G-10CR intercepted suspension.
Direct Vespel support is rejected and is not part of the current design.}

Committed elements \designspec{}:
\begin{itemize}[noitemsep]
  \item Kevlar~49 tension members ($\kappa_\mathrm{Kevlar} \approx 5\times10^{-4}$\,W/m/K at 10\,mK~\cite{Barucci2005}) \designspec{}
  \item G-10CR fiberglass flexures for lateral constraint \designspec{}
  \item Staged OFHC copper thermal intercepts at 4\,K, 1\,K, 100\,mK, 10\,mK \designspec{}
  \item Gold-plated copper anchor blocks \designspec{}
  \item Ag-sinter thermal interface to sample stage \designspec{} (45\,cm$^2$ \assumed{})
\end{itemize}

Estimated support heat load: $P_\mathrm{support} \approx 0.167$\,nW
(Kevlar/G-10CR intercepted; \assumed{} --- step-response not performed).

Gate status: \conditional{} (E05, E06; hardware not fabricated).

% =============================================================================
\section{Mode-Local Secondary Thermal Feedback (Mode~D Only)}
\label{sec:thermal}
% =============================================================================

Within Mode~D sensing, the sample temperature obeys the iterative relation:
\begin{equation}
  T_s = T_\mathrm{fridge} + \frac{P_\mathrm{total}(T_s)}{G_\mathrm{eff}}
\end{equation}
where $G_\mathrm{eff}$ \assumed{} $= 10^{-5}$\,W/K (Kapitza model~\cite{SwartzPohl1989,Pobell2007};
Ag sinter not fabricated) and $P_\mathrm{total}$ includes all Mode~D thermal loads:
conduction through wiring ($P_\mathrm{cond}$~\assumed{}), optical average load
($P_\mathrm{opt}$~\assumed{}), microwave dissipation ($P_\mathrm{mw}$~\assumed{}),
support conduction ($P_\mathrm{support}$~\assumed{}), vibration ($P_\mathrm{vib}$~\mfrspec{}),
and radiation ($P_\mathrm{rad}$).

At design parameters: $T_\mathrm{fridge} = 10$\,mK \mfrspec{},
$P_\mathrm{total} = 4131$\,pW \assumed{},
$T_s = 10.413$\,mK \assumed{}.
Margin $\approx 4\times10^4$ \assumed{} --- requires $G_\mathrm{eff}$ verification.

Individual Mode~D load terms may exhibit mode-local secondary feedback
(a higher $T_s$ increases $P_\mathrm{rad}$ and $P_\mathrm{He}$, closing a small perturbative loop).
This is evaluated within Mode~D only.
\textbf{No cross-mode thermal coupling between Mode~B LCVD loads and Mode~D loads occurs;
LCVD and sensing are temporally separated.}

All thermal budget gates (\textsc{D3}, \textsc{D4}, \textsc{D5}): \conditional{}
--- $G_\mathrm{eff}$ and all load terms are \assumed{}.

% =============================================================================
\section{Vacuum, Residual Gas, and Suppression Strategy}
\label{sec:vacuum}
% =============================================================================

\textbf{Residual gas suppression is a required implementation condition, not a demonstrated result.}

Current state: bakeout not executed; NEG not installed; cryotrap not installed.
Gates D10a, D10b, E01, E04: \blocked{}.

Required implementation (all \designspec{}):
\begin{enumerate}[noitemsep]
  \item 250\textdegree{}C / 48\,h UHV bakeout of chamber and lines~\cite{O_Hanlon2003,Jousten2016}
  \item SAES St707 NEG pump activation post-bakeout~\cite{DalPalmiero2004}
  \item 77\,K charcoal cryotrap~\cite{CERN2020}
  \item RGA all-species verification (FC-corrected for sensing-surface pressure)~\cite{Giorgi2010}
\end{enumerate}

Pre-bakeout $P_\mathrm{H_2} \approx 10^{-10}$\,Pa $\Rightarrow$
$\theta_\mathrm{H_2} \approx 0.32\%$ (above 0.1\% threshold;
\blocked{} until bakeout executed).
Post-bakeout forecast: $\theta_\mathrm{H_2} \approx 0.003\%$ (\conditional{};
\assumed{} --- not measured).

% =============================================================================
\section{Independent Gas Line Architecture}
\label{sec:gaslines}
% =============================================================================

Required isolation architecture \designspec{}:
\begin{itemize}[noitemsep]
  \item Independent $^3$He sensing line (dedicated; no dead volume shared with CH$_4$)
  \item Independent $^4$He reference line (dedicated control)
  \item CH$_4$ / LCVD precursor line (Mode~B only; isolated during Mode~D)
  \item Pumpout / purge line (Mode~C operation)
  \item Line-specific shutoff valves, pressure gauges, and cryogenic-compatible hardware
  \item No backstreaming path from CH$_4$ to helium sensing lines
\end{itemize}

Interlock IL-02: CH$_4$ and He-3 valves may not be open concurrently.
All gas lines: \designspec{}. Gate status: \conditional{} (E07, A4).

PMB / shuttered cryo-baffle for helium dosing: \designspec{}; not installed.
Provides temporal control and isotope-contrast separation.
Status of PMB: \conditional{} (E08).

% =============================================================================
\section{QCM Coverage Proxy}
\label{sec:qcm}
% =============================================================================

A 5\,MHz AT-cut QCM adjacent to the diamond surface is proposed as a coverage proxy.
Sauerbrey sensitivity~\cite{Sauerbrey1959}:
$C_f = 2f_0^2/(\rho_q v_q) = 5.66\times10^5$\,Hz/(g/m$^2$) = 0.057\,Hz/(ng/cm$^2$).
He-3 monolayer ($1.65$\,ng/cm$^2$): $\Delta f \approx 0.09$\,Hz.
CH$_4$ monolayer (estimated): $\Delta f \approx 0.15$\,Hz.

\textbf{QCM-based coverage inference remains conditional until its frequency stability
and dose correlation are demonstrated under the actual cryogenic operating conditions.}

Requirements: sub-0.1\,Hz stability; vibration isolation; thermal drift separation;
cryogenic calibration.
Status: \designspec{}; not installed.
Gate B5: \conditional{}.

% =============================================================================
\section{NV Helium Isotope Detection Model}
\label{sec:detection}
% =============================================================================

NV centers in Mode~D respond to the stochastic magnetic field produced by
diffusing $^3$He nuclear spins~\cite{KolkowitzScience2015,MazeNature2008}.
The decoherence rate contribution is~\cite{DegenRMP2017}:
\begin{equation}
  \Gamma_\mathrm{DC} = \left(\frac{\mu_0}{4\pi}\right)^2
  \frac{\gamma_\mathrm{NV}^2\,\gamma_{^3\mathrm{He}}^2\,\hbar^2\,n_s\,\tau_c}{d^6}
  \propto \frac{S(\omega{=}0,\,\tau_c)}{d^6}
\end{equation}
where $S(\omega,\tau_c) = 2\Gamma_0\tau_c/(1+\omega^2\tau_c^2)$ is the Lorentzian
spectral density of the surface magnetic noise,
$n_s$~\assumed{} is the areal spin density,
$d$~\assumed{} is the NV--surface separation, and
$\tau_c$~\unknown{} is the $^3$He surface correlation time.

Measurement protocol uses $^4$He as a non-magnetic reference:
\begin{equation}
  \Delta\Gamma = \Gamma_{^3\mathrm{He}} - \Gamma_{^4\mathrm{He}}
\end{equation}
Detection requires $\mathrm{SNR} = \Delta\Gamma / \delta\Gamma \geq 5$~\cite{TaylorNatPhys2008},
where $\delta\Gamma \propto (C_\mathrm{eff}\,T_{2,\mathrm{eff}}^*\,\sqrt{N})^{-1}$.

\subsection{$\tau_c$ and $C_\mathrm{contr}$: Co-Equal Scientific Bottlenecks}

\textbf{Detection feasibility is jointly gated by $\tau_c$ and $C_\mathrm{contr}$
at 10\,mK.}

\paragraph{$\tau_c$ [\unknown{}]}
The correlation time of $^3$He adsorbed on F-terminated diamond at 10\,mK is
completely unknown.
Graphite-surface analogs~\cite{Bernu2006,MorhardPRB1996,Greywall1993} provide
plausibility only; surface termination fundamentally affects dynamics.
Canonical threshold (v3.3; combined SNR$\geq5$ + $\varepsilon_\mathrm{thermo}<1\%$
including pulse dephasing):
\begin{equation}
  \tau_c \geq 292\,\mu\mathrm{s}
\end{equation}
Older threshold 27.7\,$\mu$s (v3.0; no pulse dephasing): \superseded{}.
See Appendix~\ref{app:tau_superseded}.

\paragraph{$C_\mathrm{contr}$ at 10\,mK [\unknown{}]}
NV ODMR contrast at 10\,mK under pulsed 532\,nm excitation is unknown~\cite{Sangtawesin2019,TrusheimPRL2020}.
Room-temperature values (typically 0.05--0.15) may not hold at 10\,mK due to
charge-state instability under cryogenic laser irradiation~\cite{RondinRMP2014}.

\textbf{If $C_\mathrm{contr} = 0$ or too low at 10\,mK, detection fails
regardless of $\tau_c$.}

Gate D12 [\unknown{}]: $C_\mathrm{contr}$ at 10\,mK not measurable without
ODMR at operating temperature.
Gate D13 [\conditional{}]: detection SNR requires $\tau_c \geq 292\,\mu$s \textbf{and}
$C_\mathrm{contr} > 0$.

% =============================================================================
\section{NV Sample Requirements}
\label{sec:nv}
% =============================================================================

All NV parameters are \assumed{} or \unknown{} unless measured on the actual sample:

\begin{center}
\small
\begin{tabular}{llll}
\toprule
Parameter & Value & Status & Required measurement \\
\midrule
$d_\mathrm{NV}$ & 10\,nm & \assumed{} & SIMS or PL depth \\
$T_2^*$ & 10\,$\mu$s & \assumed{} & Ramsey in cryostat \\
$C_\mathrm{contr}$ at 10\,mK & \unknown{} & \unknown{} & ODMR at 10\,mK \\
$\eta_\mathrm{col}$ & 0.0635 & \assumed{} & In-cryostat calibration \\
$n_s$ (surface spins) & $3.3\times10^{18}$\,m$^{-2}$ & \assumed{} & TPD or QCM on F-diamond \\
Surface termination & F-terminated & \designspec{} & XPS + AFM post-termination \\
\bottomrule
\end{tabular}
\end{center}

% =============================================================================
\section{Optical and Femtosecond Laser Constraints (Mode~B and Mode~D)}
\label{sec:optical}
% =============================================================================

Two distinct optical uses:
(1) Mode~B femtosecond laser for LCVD (proposed; \designspec{});
(2) Mode~D CW or pulsed laser for NV readout (proposed; \designspec{}).
These are temporally separated.

LCVD femtosecond laser parameters \assumed{}: $E_\mathrm{pulse}=10$\,fJ,
$f_\mathrm{rep}=1$\,MHz, $\tau_\mathrm{pulse}=100$\,fs, $r=361$\,nm,
$\eta_\mathrm{abs}=0.05$ \assumed{}.
Fluence: $F \approx 2.4\times10^{-2}$\,J/m$^2$ $\ll$ damage threshold
$\sim10^4$\,J/m$^2$~\cite{Ozkan2012,Shinoda2009} \litbound{}.

\textbf{Low average power does not by itself validate femtosecond-pulse safety;
transient fluence, nonlinear response, absorption, and recovery must be measured.}

Pulse gating: AOM switching time ($\sim$100\,ns--1\,$\mu$s) cannot gate
100\,fs pulse envelopes.
If individual pulse gating is required, a Pockels cell or equivalent ultrafast
shutter is needed.
Status: \designspec{}.
Thermal recovery: $\tau_\mathrm{thermal} \approx 42$\,ns (Debye model; \assumed{})
$\ll$ 1\,ms rep period.

Gate D7 [\conditional{}]: transient heating conditional on $\eta_\mathrm{abs}$ \assumed{}.

% =============================================================================
\section{Microwave Delivery and Heating}
\label{sec:mw}
% =============================================================================

Mode~D only.
CPW geometry \designspec{}; $P_\mathrm{mw}$ at sample $\approx 1$\,nW \assumed{}
(room-temperature source + estimated cryogenic attenuation; not measured at base).
Heating estimate: $\Delta T_\mathrm{mw} \approx 100\,\mu$K \assumed{}.
Threshold: $<104\,\mu$K (1\% of $T_s$).

\textbf{Microwave and vibration loads remain conditional until measured at the
sample/NV position under Mode~D operating conditions.}

Gate D16 [\conditional{}]: $\Delta T_\mathrm{mw}$ marginal; both $P_\mathrm{mw}$
and $G_\mathrm{eff}$ are \assumed{}.

% =============================================================================
\section{Vibration Constraints}
\label{sec:vib}
% =============================================================================

$S_\mathrm{vib}$ at NV position: $10^{-10}$\,m$^2$/Hz \mfrspec{}
(Oxford Triton specification; not measured at NV location in this setup).
Vibration-to-heat conversion requires mode-local model within Mode~D.

Gate D17 [\conditional{}]: Helmholtz geometry \designspec{}; $S_\mathrm{vib}$ not measured in situ.

% =============================================================================
\section{Helium Molecular Flow (Mode~D)}
\label{sec:flow}
% =============================================================================

$^3$He transport in the dosing geometry is in the molecular-flow (Knudsen) regime.
At dosing conditions ($T=4$\,K, $P=1\,\mu$Pa):
$\lambda_\mathrm{He} \approx 0.46$\,m $\gg L_\mathrm{char} = 10$\,mm,
giving $\mathrm{Kn} \approx 46$~\cite{Bird1994,Roth1990}.

Gate D9 status: \derivedcheck{} --- gas-kinetic check is self-consistent;
confirmed by formula from standard references; \textbf{not} a physical system
\textsc{Pass} (would require actual He-3 flow measurement in the built system).

% =============================================================================
\section{Mode~B Cryogenic LCVD / Surface-Processing Feasibility Limits}
\label{sec:lcvd}
% =============================================================================

Mode~B is a proposed pulsed surface-processing mode, not a demonstrated capability.
It is mutually exclusive with Mode~D sensing by hardware interlock (IL-01, IL-02, IL-14).
Helium is forbidden in Mode~B (Gate A14, IL-14).

\subsection*{Gas-phase truth}
Normal LCVD requires a gaseous or beam-delivered precursor at the reaction site.
At millikelvin substrate and shield temperatures, methane and other hydrocarbon
precursors will \emph{cryocondense} on cold surfaces unless delivered through a
controlled warm/differentially pumped molecular beam path. Therefore this work
\textbf{does not claim conventional chamber-fill LCVD at 10\,mK}.
Mode~B is a proposed pulsed, local, beam-fed surface-processing mode whose
deposition yield per pulse must be measured.

\subsection*{Mode~B laser configuration (separate from Mode~D NV readout)}
Mode~B LCVD laser parameters are \emph{not} the same as Mode~D NV readout
parameters. The package configuration vectors are kept separate:
\begin{itemize}[noitemsep]
  \item \textbf{Mode~B LCVD optics}: $E_{\mathrm{pulse}}^A$ \designspec{},
        $f_{\mathrm{rep}}^A$ \designspec{}, $\tau_{\mathrm{pulse}}^A$ \designspec{},
        spot radius \designspec{}, $\eta_{\mathrm{abs}}^A$ \designspec{},
        precursor species \designspec{},
        deposition yield per pulse \unknown{},
        local surface temperature during deposition \unknown{}.
  \item \textbf{Mode~D NV readout}: $E_{\mathrm{pulse}}^D$ \designspec{},
        $f_{\mathrm{rep}}^D$ \designspec{}, $\eta_{\mathrm{abs}}^D$ \assumed{},
        $T_{\mathrm{peak}}^D$ \assumed{}, $\tau_{\mathrm{recovery}}^D$ \assumed{}.
\end{itemize}
Mode~D gates D6/D7 use only the Mode~D vector. Mode~B gates A6--A13 use only
the Mode~B vector. The two vectors are never mixed in a single gate.

\subsection*{Mode~B LCVD feasibility gates}
\begin{itemize}[noitemsep]
  \item[A6] Precursor remains beam-delivered/gaseous until reaching the target;
        does not act as a chamber-filling gas. \conditional{} until measured.
  \item[A7] Precursor cryocondensation on cold shields and on the sensing zone
        is below limit. \conditional{} until measured by witness coupon and RGA.
  \item[A8] Local surface temperature during the deposition pulse is
        \emph{measured}, not assumed; photolysis or pyrolysis threshold is
        verified at that surface temperature. \blocked{} until pump--probe or
        thermometric measurement is available.
  \item[A9] Deposition yield per pulse is \emph{measured} by witness coupon /
        QCM / AFM / Raman / XPS. \blocked{} until measured.
  \item[A10] Growth-zone contamination does not reach sensing-zone coupon.
        \conditional{} until measured by witness coupon and RGA after each
        Mode~B pulse train.
  \item[A11] Byproducts and unreacted precursor pump below RGA thresholds
        \emph{before} Mode~D begins. \blocked{} until RGA is installed and
        all-species threshold is verified.
  \item[A12] Mode~B heat is removed by a dedicated non-MC thermal dump path
        (Section~\ref{sec:mode_b_thermal_dump}). \conditional{} until the
        dedicated dump switch and dump-stage temperature sensors are installed
        and the heat balance is measured.
  \item[A13] Repeated mode cycling (Mode~B processing $\rightarrow$ Mode~C recovery $\rightarrow$ Mode~D sensing) passes fatigue and stress inspection
        (Section~\ref{sec:cycling}). \conditional{} until $N=100$ qualification
        cycles are completed.
  \item[A14] He-3 / He-4 film absent before Mode~B. \blocked{} until verified
        in-situ by RGA, QCM, and pressure logging. Hardware interlock IL-14:
        \textsc{LCVD\_on} $\wedge$ \textsc{helium\_film\_present} = IMPOSSIBLE.
        Rationale: adsorbed or superfluid helium can wet surfaces, block
        precursor access, alter thermal transport, alter NV optical/magnetic
        environment, and corrupt surface chemistry. Helium belongs only to
        Mode~D sensing.
\end{itemize}
None of A6--A14 may reach \textsc{Pass} from assumptions, simulation
convergence, or literature; each requires measured in-system data.

\subsection{Mode~B non-MC thermal dump architecture}
\label{sec:mode_b_thermal_dump}
The mixing chamber (MC) cannot remove Mode~B laser heat in real time. The
proposed Mode~B thermal dump architecture is:
\begin{itemize}[noitemsep]
  \item \textbf{MC-protection switch}: superconducting heat switch held
        \emph{open} during Mode~B so the MC is decoupled from the growth zone.
  \item \textbf{Growth thermal dump switch}: an independent heat switch that
        \emph{closes} during Mode~B and connects the growth-zone carrier to a
        1\,K or 4\,K dump stage (separate from the MC).
  \item \textbf{Radiation shutter stack}: closed during Mode~B to protect the
        10\,mK region from optical and IR radiation from the growth zone.
  \item \textbf{Temperature sensors}: independent calibrated thermometers on the
        growth-zone surface, sample bulk, MC, 1\,K stage, and 4\,K stage,
        logged continuously during Mode~B.
\end{itemize}
The minimal Mode~B thermal solve is
$T_{\mathrm{growth\_surface}}(t),\;T_{\mathrm{sample\_bulk}}(t),\;T_{\mathrm{MC}}(t)$
driven by $P_{\mathrm{laser}}^A(t)$ and $P_{\mathrm{precursor}}^A(t)$ through
$G_{\mathrm{growth\_to\_dump}}$ with parasitic $G_{\mathrm{growth\_to\_MC\_leak}}$.
The requirement is that the MC heat leak stays below the MC cooling budget
($\ll 200\,\mu$W) while the growth-zone heat is absorbed by the non-MC dump.
This package does not assert that the 10\,mK MC removes LCVD heat in real time.
All conductances $G_{\cdot}$ are \designspec{} and must be measured before
Mode~B is run.

\subsection{Thermal cycling, stress, and fatigue}
\label{sec:cycling}
Repeated A/B/C/D cycling stresses material interfaces. The package requires
measurement (or, where unavailable, \conditional{} status) for:
diamond / Ag-sinter CTE mismatch;
Ag sinter cracking or delamination;
Kevlar end-termination slip after repeated cycles;
G-10CR flexure fatigue;
copper anchor stress;
wire and CPW fatigue at mode-switching cadence;
shutter actuator vibration and wear;
gas-line and feedthrough cycling.

Acceptance criteria:
no cracks; no delamination; no conductance drift greater than threshold;
no RGA contamination increase; no measurable alignment drift,
after $N=100$ cycles (first engineering qualification) and
$N=1000$ cycles (extended qualification). Status remains \conditional{} until
the package is physically cycled.

\subsection*{Required Mode~B validation evidence}
Witness coupon + QCM + AFM + Raman + XPS + RGA records must be archived for
\emph{every} Mode~B run. Mode~D is permitted to begin only after the witness
coupon, RGA, and thermometric records demonstrate that A11 (byproduct removal),
A10 (cross-contamination), A14 (helium-film absence), and A12 (thermal recovery)
are satisfied.

\textbf{Final Mode~B status: BLOCKED (hardware not installed) /
CONDITIONAL post-installation. LCVD at 10\,mK is not validated.}

\subsection{Speculative material extensions (future work only)}
The following are \textbf{not} part of the current gate verdict and require
independent experimental validation before being incorporated:
isotopically pure C-13 diamond growth, graphene-on-diamond, CNT forests,
graphene--diamond--CNT hybrids, or any post-silicon device concept.



% =============================================================================
\section{Simulation Architecture}
\label{sec:sim}
% =============================================================================

\texttt{qta\_full\_sim.py} implements a single-mode state vector
(\texttt{ModeStateVector}) that computes, in order:
(1) surface coverage estimates ($\theta_\mathrm{H_2}$, helium dose);
(2) molecular transport ($\lambda_\mathrm{He}$, $\mathrm{Kn}$);
(3) iterative thermal solve ($T_s = T_f + P(T_s)/G_\mathrm{eff}$, converged to $<1\,\mu$K);
(4) optical transient ($T_\mathrm{peak}$, $F$, $\eta_\mathrm{abs}$ \assumed{});
(5) detection model ($\Gamma_\mathrm{DC}$, SNR, $\tau_c$ threshold); and
(6) mode-local secondary load terms.

\textbf{Simulation outputs are conditional forecasts.
Numerical convergence does not imply physical validation.
No gate reaches Pass from assumed inputs.}

Run: \verb|python3 qta_full_sim.py|
Outputs written to \verb|outputs/| directory.

\subsection{Non-lumped 1D/2D reduced-order multiphysics layer}
\label{sec:multiphysics}

In addition to the lumped estimates above, QTA now includes a serious 1D/2D
\emph{non-lumped reduced-order multiphysics forecast layer}
(\texttt{qta\_multiphysics/}). It provides spatially-resolved forecast backends
for thermal transport, moving-boundary laser absorption, gas/contamination
transport, surface coverage, optical deposition, microwave heating, radiation
leakage, vibration transfer, and the coupled Mode~B $\rightarrow$ Mode~C
$\rightarrow$ Mode~D recovery cycle. The 1D moving-boundary heat equation
$\rho C_p(T)\,\partial_t T = \partial_z[k(T)\partial_z T] + Q_\mathrm{laser} +
Q_\mathrm{mw} + Q_\mathrm{bg}$, with a Kapitza-type backside sink
$-k\,\partial_z T|_L = \alpha_K(T_L^4 - T_\mathrm{fridge}^4)$, is solved by
finite-volume method-of-lines with a stiff (BDF) integrator; the 2D axisymmetric
backend solves the same physics in $(r,z)$ with cylindrical $2\pi r$ weighting, a
symmetry axis at $r=0$, and a Gaussian radial beam profile.

\textbf{1D is the canonical, fast, default non-lumped path; 2D axisymmetric is the
serious spatial refinement; 3D is future work and is neither implemented nor
claimed in this work.} The layer is numerically self-checked (mesh convergence,
source-energy conservation, Kapitza-sign, axis non-singularity, and a verified
2D$\rightarrow$1D reduction to $\sim$1.4\%); these are numerical
self-consistency checks only. Every output of the layer is model-only and carries
\texttt{measured\_in\_this\_system = false}; the legacy lumped model is retained
solely as a comparator. The layer contributes 20 additional gates, all
\conditional{}/\blocked{}/\derivedcheck{} --- \textbf{none reaches Pass}. As with
the lumped simulation, numerical convergence does not imply physical validation,
and all multiphysics outputs remain forecast-only until hardware validation.

% =============================================================================
\section{Gate Table and Current Verdict}
\label{sec:gates}
% =============================================================================

Gate status is defined by the strictest supporting evidence.
Full table in \texttt{results\_gate\_table.csv}.
Selected critical gates:

\begin{center}
\small
\begin{longtable}{lllp{6cm}}
\toprule
Gate & Status & Canonical value & Blocking reason \\
\midrule
D9 & \derivedcheck{} & Kn = 46 & Gas kinetics only; not in-system verified \\
D10a & \blocked{} & Bakeout: not executed & E01: bakeout NOT executed \\
D10b & \blocked{} & $\theta_\mathrm{H_2}=0.32\%$ & Follows D10a block \\
D12 & \unknown{} & $C_\mathrm{contr}$: UNKNOWN & ODMR at 10\,mK not performed \\
D13 & \conditional{} & $\tau_c$: UNKNOWN & Requires $\tau_c\geq292\,\mu$s AND $C_\mathrm{contr}>0$ \\
D5 & \conditional{} & $G_\mathrm{eff}$: ASSUMED & Ag sinter not fabricated \\
E01 & \blocked{} & Bakeout: NOT INSTALLED & Must execute before D10a \\
E04 & \blocked{} & RGA: NOT PERFORMED & Must verify residual gas \\
E11--E14 & \blocked{} & Various calibrations & Not performed \\
A1--A5 & \conditional{} & Various & Hardware \designspec{} \\
\midrule
\multicolumn{4}{l}{\textbf{Summary:} 0~Pass $|$ 39~\conditional{} $|$ 0~Fail
$|$ 2~\unknown{} $|$ 21~\blocked{} $|$ 1~\derivedcheck{}} \\
\bottomrule
\end{longtable}
\end{center}

\verdictbox{CURRENT SYSTEM: BLOCKED. Mode B prerequisites not executed.
Full-cycle MC = 0.0\% (logical consequence of 23 blocked gates, not a statistical result).

POST-INSTALLATION FORECAST: CONDITIONAL.
Mode D forecast MC $\approx 33.8\%$ (HYPOTHETICAL; post-installation; $\tau_c$ log-uniform $[1\,\mathrm{ns}, 100\,\mathrm{ms}]$).
Dominant failure: $\tau_c < 292\,\mu$s (66.2\% of failed samples).
CO-EQUAL BOTTLENECK: $C_\mathrm{contr}$ at 10\,mK \unknown{}.

VALIDATED SYSTEM: NOT AVAILABLE.}

% =============================================================================
\section{First Decisive Experiment}
\label{sec:firstexp}
% =============================================================================

\textbf{The first decisive experiment is not a full QTA demonstration.}

\textbf{Question:} Can shallow NV centers at 10\,mK retain usable contrast
and distinguish $^3$He from $^4$He surface-induced decoherence under controlled exposure?

Required measurements:
\begin{enumerate}[noitemsep]
  \item $C_\mathrm{contr}$ at 10\,mK via ODMR before helium (determines if sensing is possible at all)
  \item $T_2^*$ at 10\,mK via Ramsey
  \item $^4$He exposure response (non-magnetic reference)
  \item $^3$He matched exposure response
  \item Isotope contrast $\Delta\Gamma = \Gamma_{^3\mathrm{He}} - \Gamma_{^4\mathrm{He}}$
  \item Pressure/exposure-time sweep for $\tau_c$ bound
  \item Sample temperature stability and residual gas baseline
\end{enumerate}

Pass/fail logic:
If $C_\mathrm{contr} = 0$: sensing fails; stop.
If $^3$He and $^4$He both degrade equally: non-magnetic surface effect dominates.
If $^3$He uniquely increases decoherence above noise floor: candidate nuclear-spin signal.
If $\tau_c < 292\,\mu$s from any extraction: detection forecast fails.

This experiment determines whether the sensing premise survives.
Do not claim LCVD breakthrough before this measurement is performed.

% =============================================================================
\section{Limitations}
\label{sec:limitations}
% =============================================================================

\begin{enumerate}[noitemsep]
  \item No physical measurement has been performed in this system.
  \item All parameters are \assumed{}, \unknown{}, or from \mfrspec{}/\litbound{}.
  \item Mode separation is specified; interlocks are not installed or tested.
  \item $\tau_c$ on F-diamond at 10\,mK is completely unknown.
  \item $C_\mathrm{contr}$ at 10\,mK is completely unknown.
  \item $G_\mathrm{eff}$ is model-based; Ag sinter not fabricated.
  \item Bibliography contains 23 entries but full-depth citation coverage
        of every numerical parameter requires further expansion.
  \item LCVD at 10\,mK is not demonstrated.
  \item Kevlar/G-10CR support architecture is \designspec{} and not fabricated.
  \item QCM sub-Hz stability under cryogenic conditions is unverified.
\end{enumerate}

% =============================================================================
\section{Conclusion}
\label{sec:conclusion}
% =============================================================================

This work does not demonstrate a working millikelvin LCVD/NV platform.
It defines a proposed, mode-separated, same-chamber validation framework
and the experimental program required to test it.

The current system is \textbf{blocked} because required hardware is not
installed and required measurements are not verified.

Detection feasibility is governed by two co-equal scientific unknowns:
$\tau_c \geq 292\,\mu$s (canonical v3.3) and $C_\mathrm{contr} > 0$ at 10\,mK.
Neither is measured.
Either alone is sufficient to prevent detection.

Engineering feasibility remains conditional on measured $G_\mathrm{eff}$,
residual gas suppression, $\eta_\mathrm{abs}$, $P_\mathrm{mw}$, $S_\mathrm{vib}$,
QCM stability, $d_\mathrm{NV}$, $n_s$, surface termination, and interlock behavior.

Until those measurements are performed, this framework remains an \textbf{unvalidated
conditional pre-digital-twin architecture} --- not a demonstrated breakthrough.

\textbf{Global status:}
Current: \blocked{}.
Forecast (post-installation, Mode~D): \conditional{} (33.8\% hypothetical MC).
Validated: not available.
Breakthrough claim: not made.

% =============================================================================
% APPENDICES
% =============================================================================
\appendix

\section{Superseded $\tau_c$ Thresholds}
\label{app:tau_superseded}

\paragraph{v3.0 threshold: $\tau_c \geq 27.7\,\mu$s \superseded{}}
The 27.7\,\textbf{[SUPERSEDED v3.0]}\,$\mu$s value was derived for DC Ramsey sensing without pulse dephasing.
It is superseded by the v3.3 combined threshold.

\paragraph{Canonical v3.3 threshold: $\tau_c \geq 292\,\mu$s}
Derived from: SNR$\geq5$ plus $\varepsilon_\mathrm{thermo}<1\%$,
including pulse dephasing correction ($\tau_{\pi/2}/T_2^* = 1.13$; degraded regime).
This is the value used in all gates, the final verdict, and the Monte Carlo.

\paragraph{v3.1 DC Ramsey threshold: $\tau_c \geq 27.7\,\mu$s \superseded{}}
Historical derivation retained for traceability.
Not current gate logic.

\paragraph{Superseded AC model: $\tau_c \approx 65$\,ns \superseded{}}
Used in earlier simulation versions.
Not current.

\section{Rejected Vespel Baseline Calculations}
\label{app:vespel_rejected}

Direct Vespel SP-22 supports are rejected from the committed architecture.
For reference: three $\phi3$\,mm, $L=50$\,mm rods would give
$P_\mathrm{support} \approx 1.2$\,nW (no thermal interception),
creating an unacceptable direct 4\,K-to-10\,mK conductive path.
This is not the committed design.

\section{Hardware Registry}
\label{app:hardware}
See \texttt{hardware\_registry.json}.
All items: \designspec{} or \notinstalled{} unless marked otherwise.
No item is \textsc{Verified}.

\section{Validation Matrix}
\label{app:valmat}
See \texttt{validation\_matrix.csv} (157 rows; all \texttt{can\_PASS\_now = NO}).

\section{Source and Citation Audit}
\label{app:sources}
See \texttt{representative\_source\_audit.csv} and \texttt{source\_audit\_status.txt}. The included audit is \textsc{Representative\_Only} and not submission-complete. A full submission-grade numerical source audit remains incomplete.

\section{Simulation Output Schema}
\label{app:sim}
Run \verb|python3 qta_full_sim.py| to regenerate all CSV/JSON outputs.

\section{Slop Removal Log}
\label{app:slop}
See \texttt{slop\_removal\_log.txt} for before/after record of removed
overconfidence language.

% =============================================================================
\section*{Bibliography}
% =============================================================================
\begin{thebibliography}{99}

\bibitem{Pobell2007} F. Pobell, \textit{Matter and Methods at Low Temperatures}, 3rd ed.\ Springer (2007).
\bibitem{SwartzPohl1989} E.T. Swartz and R.O. Pohl, Rev.\ Mod.\ Phys.\ \textbf{61}, 605 (1989).
\bibitem{Barucci2005} M. Barucci \emph{et~al.}, Cryogenics \textbf{45}, 295 (2005).
\bibitem{DegenRMP2017} C.L. Degen, F. Reinhard, and P. Cappellaro, Rev.\ Mod.\ Phys.\ \textbf{89}, 035002 (2017).
\bibitem{RondinRMP2014} L. Rondin \emph{et~al.}, Rep.\ Prog.\ Phys.\ \textbf{77}, 056503 (2014).
\bibitem{MazeNature2008} J.R. Maze \emph{et~al.}, Nature \textbf{455}, 644 (2008).
\bibitem{KolkowitzScience2015} S. Kolkowitz \emph{et~al.}, Science \textbf{347}, 1129 (2015).
\bibitem{TrusheimPRL2020} M.E. Trusheim \emph{et~al.}, Phys.\ Rev.\ Lett.\ \textbf{124}, 020501 (2020).
\bibitem{Sangtawesin2019} S. Sangtawesin \emph{et~al.}, Phys.\ Rev.\ X \textbf{9}, 031052 (2019).
\bibitem{TaylorNatPhys2008} J.M. Taylor \emph{et~al.}, Nat.\ Phys.\ \textbf{4}, 810 (2008).
\bibitem{Bernu2006} B. Bernu \emph{et~al.}, Phys.\ Rev.\ Lett.\ \textbf{96}, 255503 (2006).
\bibitem{MorhardPRB1996} D. Morhard \emph{et~al.}, Phys.\ Rev.\ B \textbf{53}, 14978 (1996).
\bibitem{Greywall1993} D.S. Greywall, Phys.\ Rev.\ B \textbf{47}, 309 (1993).
\bibitem{Bird1994} G.A. Bird, \textit{Molecular Gas Dynamics and the Direct Simulation of Gas Flows}, Oxford (1994).
\bibitem{Roth1990} A. Roth, \textit{Vacuum Technology}, 3rd ed., Elsevier (1990).
\bibitem{O_Hanlon2003} J.F. O'Hanlon, \textit{A User's Guide to Vacuum Technology}, 3rd ed., Wiley (2003).
\bibitem{Jousten2016} K. Jousten (ed.), \textit{Handbook of Vacuum Technology}, Wiley-VCH (2016).
\bibitem{CERN2020} CERN, \textit{Outgassing Handbook}, 2nd ed.\ (2020).
\bibitem{DalPalmiero2004} A. Dal Palmiero \emph{et~al.}, Vacuum \textbf{73}, 121 (2004).
\bibitem{Giorgi2010} D. Giorgi \emph{et~al.}, J.\ Vac.\ Sci.\ Technol.\ A \textbf{28}, 1332 (2010).

% QCM
\bibitem{Sauerbrey1959} G. Sauerbrey, Z.\ Phys.\ \textbf{155}, 206 (1959).
\bibitem{Ozkan2012} A. Ozkan \emph{et~al.}, Appl.\ Surf.\ Sci.\ \textbf{259}, 1 (2012).
\bibitem{Shinoda2009} M. Shinoda \emph{et~al.}, Appl.\ Phys.\ A \textbf{94}, 305 (2009).
\end{thebibliography}

\end{document}
